\section{Comparison: Part 1 vs. Part 2}

As established in Section 3.1, both parts of the practical project built the same application with the same underlying model, Claude Sonnet 4.6, held constant so that any differences between the two parts could be attributed to how the collaboration was structured rather than to a difference in model capability. Part 2 did not start over: it continued Part 1's codebase, inheriting its architecture, its problems, and its catalogued weaknesses. What changed was not the software being built or the model building it, but the process around it.

The persistent context document each part relied on differed in when its rules existed relative to the problems they addressed. CLAUDE.md began with only a partial set of conventions and grew reactively, each addition following a problem already encountered. AGENTS.md existed in a mature form from the outset and bound all three agent roles uniformly; several of its rules --- the non-nullable foreign key requirement is one example --- are traceable directly to a specific mistake already made in Part 1, rather than arrived at independently.

The two parts also differ in whether a design was reviewed before implementation began. Part 1 had no per-feature planning step: prompts were sequenced ad hoc, and a design was never drafted and approved ahead of writing code. Part 2 made this mandatory through \textbf{\texttt{implementation\_plan.md}}'s Design Review, which cannot be skipped or left incomplete, and which surfaced at least one concern --- the privacy implications of the super-admin role --- before any code existed to review.

This is inseparable from a shift in who did the work. Part 1 was a single agent working through one continuous thread, handling requirements interpretation, implementation, and whatever review occurred together. Part 2 split these into three roles --- Architect, Implementer, Post-Implementation Review Agent --- coordinated through structured artifact handoff rather than one undifferentiated exchange.

Human oversight moved correspondingly. In Part 1 it existed but operated after output already existed, and in practice changes were accepted once they worked rather than checked against the system they joined; oversight only became substantive once a problem was already obvious, as with MainPage.js's growth. In Part 2, human approval is a gate before implementation begins, and a vague or incomplete plan must be rejected rather than corrected later.

Testing followed the same trajectory from informal to formalized. Part 1's tests were backend-only, written after the implementation they verified, and evaluated once, after the fact, along five ad hoc dimensions; no frontend tests existed at all. Part 2 replaced this with a fixed three-tier strategy and a rule tying every test to a stated requirement rather than to whatever the implementation happens to do, so passing tests alone no longer imply a requirement was actually covered.

The handling of a failing test shows the same pattern at a smaller scale. Part 1's Test Failure Analysis Principle emerged from one specific incident --- a collections test that was made to pass by patching around a workspace-seeding inconsistency rather than fixing it --- and was applied after the fact, as a lesson drawn from that case rather than a standing rule enforced elsewhere in the project. Part 2 wrote the identical reasoning into the Charter as a mandatory step: before changing anything, an agent must read the test, read the implementation, and explicitly determine and record whether the fault lies with the implementation, the test's own expectation, or an underlying requirement that was never clearly defined --- required for any failing test, regardless of how minor the failure looks, rather than reserved for cases serious enough to prompt that reasoning on their own.

Accumulated weaknesses were also handled differently. In Part 1, each major architectural problem was only identified once it had already spread, and fixed as a retrofit; fifteen further weaknesses were catalogued and left unresolved at handoff. In Part 2, addressing an inherited weakness went through the same Architect--human--Implementer--Reviewer sequence as any new feature --- fittingly, the first feature run through that sequence was a fix to the MainPage.js monolith itself --- and a separate, independent security audit added a second detection mechanism beyond the inherited list.

Finally, documentation consistency: Part 1's OpenAPI specification had no automated check tying it to the implementation. Part 2 added a pytest module comparing the specification against the application's actual routes at runtime, folding drift detection into the same testing cycle already required for everything else.

Across all of these, the pattern is consistent: what could be built did not change, since the same model produced both parts, but where a decision was checked, by whom, and at what point in the process shifted from informal and retrospective to formalized and positioned ahead of implementation. Whether that shift produced better software, rather than more process, is taken up in the evaluation that follows.

